%%%%%%%%%%%%%%%%%%%%%%%%%%%%%%%%%%%%%%%%%%%%%%%%%%%%

\section{Ontologias}
\label{sec:ontologia}

O eProfile utiliza duas ontologias: uma ontologia de trilhas e uma ontologia de perfis. As ontologias utilizadas pelo eProfile são ontologias básicas, estendidas através de subclasses de acordo com o formato da trilha enviada pelo gerenciador (no caso da ontologia de trilha) ou das necessidades da aplicação (no caso da ontologia do perfil). Isto é exemplificado na Figura~\ref{fig:heranca}, onde o eProfile oferece uma classe genérica ``Atividade'', que é especializada pela classe ``AtividadeDeCompra'', utilizada por um dos gerenciadores de trilha.

\begin{figure}
  \centering\includegraphics[width=0.6\textwidth]{imagens/heranca.png}
  \caption{Exemplo da extensão de uma classe do eProfile}
  \label{fig:heranca}
\end{figure}

%%%%%%%%%%%%%%%%%%%%%%%%%%%%%%%%%%%%%%%%%%%%%%%%%%%%

\subsection{Ontologia de Trilha (eOntoTrail)}
\label{sec:eontotrail}

O eProfile implementa uma ontologia de trilha, chamada \textbf{eOntoTrail}. Esta ontologia será alimentada a partir dos diferentes gerenciadores de trilhas suportados, onde o agente conversor é responsável por receber os dados no formato do gerenciador de trilhas e convertê-los para a eOntoTrail. Assim, o raciocinador terá uma ontologia única sobre a qual poderá raciocinar.

A eOntoTrail está representada na Figura~\ref{fig:ont_trilha}, onde todos os relacionamentos, representados pelas setas, são do tipo ou ``has-a'' ou ``has-some''.

%\todo[inline]{Está correto utilizar ``has-a'' e ``has-some''? Estou querendo representar ``0 ou 1'', e ``0 ou mais'' respectivamente.}

\begin{figure*}
  \centering\includegraphics[width=1\textwidth]{imagens/ont_trilha_simples.png}
  \caption{Ontologia de Trilha do eProfile}
  \label{fig:ont_trilha}
\end{figure*}

\citetexto{abowd99} definem que aplicações sensíveis a contexto estão interessadas em saber {\bf quem}, {\bf onde}, {\bf quando} e {\bf o quê} (isto é, o que o usuário está fazendo), de forma a poderem descobrir por que a situação está ocorrendo. Assim, uma trilha é representada por um conjunto de eventos caracterizados pelas seguintes classes:

\begin{itemize}

  \item {\em Entity:} é a entidade que está realizando a ação. Esta entidade é vinculada à ontologia descrita na Seção~\ref{sec:eontoprofile};

  \item {\em Participants:} uma lista de outras entidades relacionadas à ação. No caso de uma reunião de negócios, por exemplo, o perfil que está sendo analisado é do tipo {\em Entity}, enquanto os outros participantes são membros da classe {\em Participants};

  \item {\em Location:} a localização onde está ocorrendo a ação. Uma localização pode ser uma localização física específica (como uma sala) ou não (uma estrada, por exemplo). A localização pode estar vinculada a uma região que, opcionalmente, pode estar vinculada a outras subregiões e geralmente possui uma localização global (GPS) vinculada a ela. Outros tipos de localização podem ser estendidas a partir da classe;

  \item {\em Time:} data e hora em que ocorreu o evento;

  \item {\em Activity:} refere-se à ação que a entidade está executando (ou tentou executar). É a classe mais abrangente da ontologia, e a que possivelmente será a mais especializada pelos gerenciadores de trilha.

\end{itemize}

%É importante notar que a própria classe base {\em Thing}\footnote{Em OWL, a classe {\em Thing} é a classe base a partir da qual todas as outras classes são extendidas.} pode ser extendida, como forma de registrar outras informações necessárias. Por exemplo, a trilha registra o interesse de um usuário em um determinado produto, a classe {\em Thing} pode ser extendida à classe {\em Product} para o registro das informações deste produto.

%%%%%%%%%%%%%%%%%%%%%%%%%%%%%%%%%%%%%%%%%%%%%%%%%%%%

\subsection{Ontologia de Perfil (eOntoProfile)}
\label{sec:eontoprofile}

A ontologia de perfil é a ontologia utilizada para armazenar os modelos de entidade, bem como as informações de perfil que são inferidas a partir da trilha. Para que a ontologia seja viável dentro do modelo proposto, ela deve ser uma ontologia genérica que permita que suas classes sejam estendidas pelas aplicações. %Uma ontologia genérica não está limitada a um domínio específico do conhcecimento,  Existem algumas ontologias genéricas propostas, tais como FOAF \cite{foaf}, GUMO \cite{gumo}, Ontonym \cite{ontonym}, e a ontologia proposta por \citetexto{golemati07}.

A ontologia proposta, a qual chama-se eOntoProfile, está representada na Figura~\ref{fig:ont_perfil}. Nesta figura, um relacionamento entre as classes é mostrado através de uma linha pontilhada, e uma especialização de uma classe é mostrado por uma linha contínua. A eOntoProfile é uma ontologia básica, que será estendida pelas aplicações de modo a montarem seus próprios modelos de entidade.

\begin{figure}
  \centering\includegraphics[width=0.8\textwidth]{imagens/ont_perfil2.png}
  \caption{Ontologia de Perfil do eProfile}
  \label{fig:ont_perfil}
\end{figure}

A classe básica da ontologia é a classe {\em Element} (elemento). A classe {\em Element} representa qualquer coisa no mundo real. Ela tem como filhos as classes:

\begin{itemize}

  \item {\it Entity:} representa uma entidade, isto é, um elemento que tenha um perfil associado a ela. Por exemplo, se uma pessoa tem um perfil no eProfile e é registrado um evento de trilha associado a outra pessoa {\bf que também possua um perfil}, esta segunda pessoa é registrada como uma {\it Entity} neste evento;

  \item {\it NonEntity:} representa um elemento que não possui um perfil associado a ela. Por exemplo, uma trilha registra o evento em que uma pessoa entra em uma loja, mas a loja não tem um perfil no eProfile. Neste caso, a loja é uma {\it NonEntity};

\end{itemize}

Esta diferenciação é importante porque no primeiro caso, o eProfile tem acesso a uma vasta gama de informações a respeito da entidade, das entidades relacionadas àquela entidade, e assim por diante. No caso de uma {\em NonEntity}, o perfil não está disponível, e é apenas um nome para o eProfile.

A ontologia FOAF \cite{foaf} utiliza o conceito de {\em knows} (conhece) para vincular duas pessoas. A eOntoProfile estende este conceito através da propriedade {\em relates-to} (se-relaciona-com) que vincula dois elementos. Através desta propriedade, é possível não apenas vincular os elementos, mas também descrever o tipo de relação que eles tem. Exemplos de relações são {\em friend-of} (amigo-de), representando uma amizade entre duas pessoas, {\em employee-of} (empregado-de), representando um vínculo empregatício entre duas pessoas, {\em owns} (possui), representando que uma pessoa possui um dispositivo, {\em depends} (depende), onde um dispositivo depende de outro, e qualquer outra relação que a aplicação deseje criar.

Uma entidade pode conter diversas características. As características das entidades são inspiradas no modelo proposto por \citetexto{golemati07}, embora o modelo tenha sido simplificado:
 
\begin{itemize}

  \item {\em PersonalInfo:} contém quaisquer informações próprias da entidade, tais como nome, localização atual, características, entre outros;

  \item {\em Interest:} quaisquer interesses ou preferências que a entidade tenha, como preferência musical ou preferência por um tema de cores específico em suas aplicações;

  \item {\em Knowledge:} conhecimento possuído pela entidade. Pode referir-se a conhecimento formal (tal como formação acadêmica) ou informal (habilidade ou experiência em alguma atividade).

\end{itemize}

Os modelos de entidade estenderão estas categorias de acordo com as suas necessidades.

Todos os modelos de entidade definidos para um perfil serão estendidos em uma mesma ontologia. Isto possibilitará a interpersonalização, ou seja, permitirá com que os dados gerados por uma aplicação possam ser utilizadas por outra. Assim, existe a possibilidade de conflito de nomes, ou seja, uma aplicação pode criar uma classe chamada {\em Preference} onde ela registra preferências da entidade por um determinado tipo de livros, enquanto outro aplicativo cria uma classe com o mesmo nome para registrar as preferências de tela da aplicação.

Para evitar este conflito, um nome único de aplicação deve ser concatenado ao nome da classe. Assim, se um dos aplicativos utilizasse o nome {\em books} e o outro utilizasse o nome {\em gps}, teriam-se, respectivamente, as classes {\em books.Preference} e {\em gps.Preference}.

Para que um aplicativo possa ter acesso aos dados gerados por outro, existem duas possibilidades:

\begin{enumerate}

  \item um aplicativo acessar diretamente os dados gerados pelo outro. No exemplo acima, o aplicativo {\em books} acessaria diretamente os dados da classe {\em gps.Preference};

  \item os aplicativos utilizarem um contrato onde eles concordam em utilizar os mesmos campos. Assim, {\em books} e {\em gps} poderiam utilizar o nome {\em findlibrary}, criando o campo {\em findlibrary.Preference} onde ambos gravariam e leriam as informações.

\end{enumerate}

Todos as extensões (subclasses) da ontologia são definidas pelas aplicações, e todas as informações do perfil são inferidas. A forma como esta inferência ocorre é descrita na Seção~\ref{sec:regras}.
